高速公路作为国家综合交通运输体系中的关键基础设施，具有通行速度高、交通流量大、运行场景连续性强等特点。在高速公路运营过程中，团雾、平流雾、辐射雾等低能见度天气会显著削弱驾驶员的前方可视距离，导致目标识别延迟、制动决策滞后和追尾连撞风险上升。与普通城区雾天不同，高速公路团雾往往具有局地突发、空间不均匀、持续时间短但危害强的特点，因此仅依赖固定阈值管控或人工视频巡查难以满足现代高速公路安全管理对实时性、准确性和可解释性的要求。面向这一问题，本文结合道路视频监控场景与智能交通实际需求，围绕“车辆检测 + 雾类型识别 + 雾浓度回归”三项任务，设计并实现了一套基于深度估计与多任务学习的高速公路团雾监测方法。

本文首先从真实工程约束出发，针对团雾样本难获取、雾浓度标签难标注、现场部署算力有限等问题，提出一种“晴天交通视频数据复用 + 单目深度估计 + 在线物理造雾 + 统一多任务训练”的技术路线。具体而言，本文以 UA-DETRAC 交通监控数据为基础数据源，通过序列级划分构建训练集与验证集；利用 \code{MiDaS_small} 对原始清晰图像进行深度预计算并形成缓存，以降低训练阶段的重复推理开销；进一步基于大气散射模型构建 GPU 在线造雾模块，在训练过程中动态生成 clear、uniform、patchy 三类天气样本，并同步给出雾类型监督与散射系数 beta 监督。在模型设计上，本文使用 YOLO 检测骨干作为共享特征提取器，在高层语义特征图上附加雾分类头和 beta 回归头，从而实现检测、分类和回归三任务的统一建模。该设计既保留了 YOLO 在道路目标检测任务中的成熟表达能力，又能够在统一特征空间内学习雾天气与交通目标之间的耦合关系。

在系统实现层面，本文进一步完成了训练、推理、可视化和部署导出闭环。训练脚本支持 checkpoint 续训、缺失深度缓存自动补算、AMP 混合精度训练、非有限梯度监测以及 QAT/INT8 转换路径；推理系统在视频帧上执行 letterbox 预处理、三任务联合前向、beta 指数滑动平均与动态置信度调节，并将车辆检测结果与雾状态联合绘制到输出视频中。针对真实监控视频中“长得像车的静止背景物体会被持续误检为车辆”的问题，本文还新增了轨迹级时序过滤模块，对逐帧检测结果进行跨帧关联、静止候选判别和二次外观复核，并进一步打通了“误检轨迹日志导出 - 可疑 patch 挖掘 - 人工复核包生成 - bootstrap hard negative 收口 - 二次车辆 / 非车辆分类器训练”的闭环。考虑到真实视频场景中车辆检测与雾估计在阶段性训练条件下可能存在精度不均衡，本文还从工程角度实现了“YOLO11n 车辆检测 + 统一模型雾估计”的混合推理方案，并在结果视频右下角生成空间雾浓度分布图，用于增强系统演示的可解释性。

从研究意义看，本文工作的价值主要体现在三个方面：第一，在数据层面，给出了一种不依赖大规模真实团雾标注样本、能够复用现有晴天交通视频的数据构造方法；第二，在模型层面，探索了高速公路团雾监测中检测、天气识别与雾浓度估计的统一多任务建模思路；第三，在工程层面，形成了一条从数据预处理、模型训练到视频推理和部署导出的完整实现链路，可为后续开展团雾预警、可变限速联动、交通气象服务和边缘侧部署验证提供技术基础。实验与系统运行结果表明，本文方法能够实现对视频中车辆目标与雾状态的同步分析，并可为后续面向高速公路低能见度预警的研究与应用提供可复现、可扩展的工程原型。
